\documentclass[11pt,a4paper]{article}
\usepackage[margin=24mm,headheight=14pt]{geometry}
\usepackage[T1]{fontenc}
\usepackage{lmodern,microtype,amsmath,amssymb,amsthm,mathtools,mathrsfs}
\usepackage{booktabs,array,longtable,xcolor,hyperref,listings,fancyhdr}
\hypersetup{colorlinks=true,linkcolor=blue!45!black,citecolor=blue!45!black,bookmarksdepth=2,
urlcolor=blue!45!black,pdftitle={Swisher's (H.3): Ordinary and Supersingular Primes}}
\newtheorem{theorem}{Theorem}[section]
\newtheorem{lemma}[theorem]{Lemma}
\newtheorem{proposition}[theorem]{Proposition}
\newtheorem{corollary}[theorem]{Corollary}
\newcommand{\Z}{\mathbb Z}
\newcommand{\Q}{\mathbb Q}
\setlength{\emergencystretch}{2em}
\pagestyle{fancy}\fancyhf{}
\fancyhead[L]{\small Swisher's (H.3): unified proof manuscript}
\fancyhead[R]{\small Ordinary and supersingular primes}
\fancyfoot[C]{\thepage}
\lstset{language=Python,basicstyle=\ttfamily\scriptsize,breaklines=true,
columns=fullflexible,keepspaces=true,showstringspaces=false,frame=single,
rulecolor=\color{black!25},aboveskip=10pt,belowskip=10pt}
\title{Swisher's (H.3):\\ Ordinary and Supersingular Primes\\[7pt]
\large A Unified Proof Manuscript with the Small-Prime Correction}
\author{Research manuscript}
\date{Unified edition: 10 November 2026}
\begin{document}
\maketitle
\begin{abstract}
We present a unified proof of the ordinary and supersingular branches of
Swisher's conjecture (H.3) for $p>3$, together with the corrected
supersingular statement at $p=3$. The common input is an integral Hecke
trace in the raw cubic parameter. An explicit extra normalizer makes
this parameter a Hauptmodul and gives an all-level formal congruence
with precision $p^{3s}$ for every odd prime. At ordinary primes,
canonical CM specialization gives the multiplier
$-\Gamma_p(1/4)^4$, identified exactly using Guo's weaker published
recurrence; arbitrary raw lengths and the full cutoff are also covered.
At supersingular primes, the value and first two derivatives of a
Legendre pullback are transported to the full integral crystal. An
exact differential residual and a weighted Taylor estimate give the
two-step modulus $p^{3s-1}$ for $p\ge7$ and $3^{3s-2}$ for $p=3$.
The common trace argument is included once, both specialization arguments
are included in full, and independent finite identities and reproducible
verification programs are supplied as appendices.
\end{abstract}
\noindent\textit{Status.} This is a newly written research proof manuscript.
The arguments have undergone internal mathematical and exact computational
checks. External expert review and a definitive assessment of novelty
remain pending. Standard modular and crystalline inputs are cited
explicitly; the new connecting arguments are written out below.
\setcounter{tocdepth}{1}
\tableofcontents
\clearpage

\section{The combined statement and its precise scope}\label{sec:statement}

For an odd prime $p$ and integers $m\ge1$, $s\ge0$, define
\[
 A_{m,s}=\sum_{0\le2k<mp^s}\frac{(1/2)_k^3}{k!^3},
 \qquad H_s=A_{1,s},\qquad H_0=1,
 \qquad
 \epsilon_p=\begin{cases}1,&p=3,\\0,&p\ge5.\end{cases}
\]
The prime is suppressed in this notation; $H_s(3)$ will indicate the
specialization $p=3$ when needed. Here $(a)_k$ is the rising factorial,
$v_p(p)=1$, and $\Gamma_p$ denotes Morita's $p$-adic gamma function,
the continuous extension to $\mathbb Z_p$ of
\[
 \Gamma_p(n)=(-1)^n\prod_{\substack{1\le j<n\\p\nmid j}}j
 \qquad(n\ge1).
\]

\begin{theorem}[Ordinary and supersingular recurrences]\label{thm:main}
Let $p$ be an odd prime.
\begin{enumerate}
\item If $p\equiv1\pmod4$, then for every $m,s\ge1$,
\begin{equation}\label{eq:ordinarytarget}
 A_{m,s}\equiv-\Gamma_p(1/4)^4A_{m,s-1}
                  \pmod{p^{3s+3v_p(m)}}.
\end{equation}
\item If $p\equiv3\pmod4$, then $H_1\equiv0\pmod{p^2}$ and, for
every $s\ge2$,
\begin{equation}\label{eq:sstarget}
 H_s\equiv p^2H_{s-2}\pmod{p^{3s-1-\epsilon_p}}.
\end{equation}
\end{enumerate}
\end{theorem}

For $p>3$, the half-length case $m=1$ of the first assertion, the
supersingular first level, and the supersingular two-step recurrence
are the three statements listed as (H.3) in \cite[Table~2, (H.3)]{Swisher}.
The theorem treats every level in their stated ranges. Its $p=3$
conclusion uses the necessary correction discussed below.

\begin{center}
\begin{tabular}{@{}llll@{}}\toprule
Prime range & Sum and levels & Recurrence & Modulus\\\midrule
$p\equiv1\pmod4$ & $A_{m,s}$, $m,s\ge1$
 & one step & $p^{3s+3v_p(m)}$\\
$p\equiv3\pmod4$ & $H_1$ & $H_1\equiv0$ & $p^2$\\
$p\ge7$, $p\equiv3\pmod4$ & $H_s$, $s\ge2$
 & two steps & $p^{3s-1}$\\
$p=3$ & $H_s$, $s\ge2$ & two steps & $3^{3s-2}$\\\bottomrule
\end{tabular}
\end{center}

The ordinary full-cutoff case is $m=2$:
\[
 \sum_{k=0}^{p^s-1}\frac{(1/2)_k^3}{k!^3}
 \equiv-\Gamma_p(1/4)^4
 \sum_{k=0}^{p^{s-1}-1}\frac{(1/2)_k^3}{k!^3}
 \pmod{p^{3s}}.
\]
The arbitrary-length and full-cutoff extensions in this manuscript
concern ordinary primes only. The supersingular assertion concerns
the exact half cutoff $(p^s-1)/2$.

\subsection{The necessary correction at three}

The unqualified exponent $3s-1$ fails at $p=3$, already at $s=2$:
\[
 H_2(3)=\frac{2576763}{2097152},\qquad
 H_2(3)-9=-\frac{3^4\cdot201205}{2^{21}}.
\]
Thus $v_3(H_2(3)-9)=4$, whereas $3s-1=5$. The exponent in
\eqref{eq:sstarget} is $3s-2$ at $3$, and is attained by this example.
This is a correction to the small-prime statement, not an unresolved
exception in the proof of the corrected theorem.

\subsection{Common input and the two specialization arguments}

Sections~\ref{sec:rawparameter}--\ref{sec:rawtrace} establish the common
formal congruence directly. Retaining the raw parameter makes its
$p^s$ truncation exactly the half-length sum in (H.3). The argument
uses the modular framework of \cite{Beukers,Katz}, with the
excellent-lift viewpoint of \cite{BV3}, and includes its own Hecke
trace induction. No result from a separate unpublished manuscript is
assumed.

Sections~\ref{sec:ordinary}--\ref{sec:gamma} treat ordinary primes.
The point $u=1/64$ is a zero of the second Hasse--Witt factor in the
BV3 cubic chart, so specialization is performed on an ordinary
modular-curve chart. Guo proved the half- and full-cutoff recurrences
modulo $p^{s+1}$ in \cite{Guo}, equations (1.9)--(1.10) in the arXiv
version. Only the half-cutoff comparison is used to identify the
constant multiplier; the higher precision comes from the trace.

Sections~\ref{sec:jets}--\ref{sec:sscm} start afresh from the all-odd-prime
formal theorem on the Legendre base. They do not specialize an ordinary
canonical lift at a supersingular point. An integral rank-three jet
congruence is first transported to a regular lift on the full crystal,
and only then evaluated at $\lambda=1/2$. This proof does not depend
on the ordinary CM or gamma-multiplier conclusion.

\section{The raw cubic parameter and the half-length truncation}\label{sec:rawparameter}

Put
\[
 \mathcal F(X)=\sum_{k\geq0}\binom{2k}{k}^{3}X^{2k},\qquad
 \mathcal F_L(X)=\sum_{0\leq 2k<L}\binom{2k}{k}^{3}X^{2k}.
\]
Thus the sum in Swisher's conjecture (H.3) is exactly
\[
 H_s=\mathcal F_{p^s}(1/8)
     =\sum_{k=0}^{(p^s-1)/2}\frac{(1/2)_k^3}{k!^3}.
\]
This identity does not discard a tail of a longer truncation.
For $q=e^{2\pi i\tau}$ define
\[
 x(\tau)=\frac{\eta(2\tau)^{12}\eta(8\tau)^{12}}
                   {\eta(4\tau)^{24}},\qquad
 \Theta(\tau)=\frac{\eta(4\tau)^{20}}
                   {\eta(2\tau)^8\eta(8\tau)^8}.
\]
The eta products and theta identities used here are classical; see
\cite[Appendix~C, Section~9.1]{Beukers} and \cite[Section~20.7]{DLMF}.
The product expansions give $x\in q+q^2\mathbb Z[[q]]$ and
$\Theta\in1+q\mathbb Z[[q]]$.  If $z=4\tau$, the usual theta identities give
\begin{equation}\label{raw:theta}
 x=\frac{\vartheta_2(z)^2\vartheta_4(z)^2}
             {4\vartheta_3(z)^4},\qquad
 \Theta=\vartheta_3(z)^4.
\end{equation}
Writing $\lambda(z)=\vartheta_2(z)^4/\vartheta_3(z)^4$, one has
\begin{equation}\label{raw:lambda}
 16x^2=\lambda(1-\lambda),\qquad
 \mathcal F(x)=
 {}_2F_1(1/2,1/2;1;\lambda)^2=\Theta.
\end{equation}
The second identity is the classical quadratic transformation followed by
Clausen's identity, with the branch specified by the expansion at $q=0$.
We use the determinant-normalized slash action
$(f|_2B)(\tau)=\det(B)(c\tau+d)^{-2}f(B\tau)$
for $B=\left(\begin{smallmatrix}a&b\\c&d\end{smallmatrix}\right)$.

\subsection{The extra normalizer and its Hauptmodul}

Let
\[
 \Gamma_0=\Gamma_0(16),\qquad
 R=\begin{pmatrix}1&1/2\\8&5\end{pmatrix},\qquad
 W=\begin{pmatrix}0&-1\\16&0\end{pmatrix},\qquad
 G=\langle\Gamma_0,R,W\rangle.
\]
Matrices are interpreted projectively when acting on the upper half plane.
This is not merely the usual Atkin--Lehner extension
$\Gamma_0(16)+16$: the additional coset represented by $R$ is essential.
In particular one cannot apply an Atkin--Lehner-only theorem without checking
the argument for this extra normalizer.

\begin{lemma}\label{raw:group}
Both $R$ and $W$ normalize $\Gamma_0$.  The quotient $G/\Gamma_0$ has
order four.  Moreover
\[
 x\circ R=x,\quad x\circ W=x,\qquad
 \Theta|_2R=\Theta,\quad \Theta|_2W=-\Theta.
\]
The function $x$ is a Hauptmodul for $G$.  The zeros of $\Theta$ are precisely
the two $\Gamma_0$-cusps with denominator $4$, and these form one $G$-orbit.
At either of these cusps the orders of $x$ and $\Theta$ on $X_0(16)$ are
$-2$ and $2$, respectively.
\end{lemma}

\begin{proof}
The eta-quotient criterion applies at level $16$.  For $x$, the two exponent
sums are $24$ and $24$, and its character is trivial.  For $\Theta$, the two
sums are zero and its character is also trivial.  Ligozat's cusp-order formula \cite[Proposition~3.2.8]{Ligozat}
gives the following complete table.  A denominator-$c$ cusp has width
$16/\gcd(c^2,16)$.
\[
\begin{array}{c|rrrrr}
 c&1&2&4&8&16\\ \hline
 \text{number of cusps}&1&1&2&1&1\\
 \operatorname{ord}_c(x)&1&1&-2&1&1\\
 \operatorname{ord}_c(\Theta)&0&0&2&0&0\\
 \text{width}&16&4&1&1&1
\end{array}
\]
The eta products have neither zeros nor poles in the upper half plane.

For an integral matrix $B=\left(\begin{smallmatrix}a&b\\16c&d\end{smallmatrix}\right)$
with odd determinant, direct multiplication gives
\begin{equation}\label{raw:conjugation}
 R^{-1}BR=
 \begin{pmatrix}
 5a+40b-8c-4d & (5a-5d)/2+25b-4c\\
 8(d-a)-64b+16c & -4a-40b+8c+5d
 \end{pmatrix}.
\end{equation}
Here $a,d$ are odd, so all entries are integral and the lower left entry is
divisible by $16$.  This proves the required inclusion for $\Gamma_0$.
The reverse inclusion follows in the same way, or from
\[
 R^2=\begin{pmatrix}5&3\\48&29\end{pmatrix}\in\Gamma_0.
\]
The normalization by $W$ is the usual Fricke normalization.  In addition
$W^{-1}RW=R^{-1}$.  Hence $G/\Gamma_0$ has at most four elements.
The cosets represented by $1,R,W,RW$ are distinct: after scaling to determinant
one, the latter three cannot be integral matrices in $\Gamma_0$.
Thus the index is exactly four.

Conjugation by $D_4=\operatorname{diag}(4,1)$ sends $\Gamma_0$ to
$\{\pm\gamma:\gamma\in\Gamma(4)\}$, sends $W$ projectively to
$S=\left(\begin{smallmatrix}0&-1\\1&0\end{smallmatrix}\right)$, and sends $R$ to
\[
 R'=\begin{pmatrix}1&2\\2&5\end{pmatrix}
   =\begin{pmatrix}1&0\\2&1\end{pmatrix}
    \begin{pmatrix}1&2\\0&1\end{pmatrix}.
\]
For the theta expression in \eqref{raw:theta}, each of the two factors on the
right changes the sign of $x$, while their product preserves it.
The transformation $S$ interchanges $\vartheta_2$ and $\vartheta_4$, with
the common automorphy factors canceling in $x$.  The form
$\vartheta_3(z)^4$ has trivial weight-two character on $\Gamma(2)$, whereas
$\vartheta_3(-1/z)^4=-z^2\vartheta_3(z)^4$.  These facts prove the four
transformation assertions.

The divisor table shows that $x:X_0(16)\to\mathbb P^1$ has degree four.
It factors through the degree-four quotient by $G/\Gamma_0$, so the induced
map $X(G)\to\mathbb P^1$ has degree one.  This proves the Hauptmodul assertion.
Finally $W(1/4)=-1/4$, so $W$ exchanges the two denominator-$4$ cusp classes.
\end{proof}

\section{The integral Hecke trace at every raw cutoff}\label{sec:rawtrace}

For an odd prime $p$, put
\[
 M(16,p)=\left\{\begin{pmatrix}a&b\\16c&d\end{pmatrix}:
 a,b,c,d\in\mathbb Z,\ ad-16bc=p\right\}.
\]
Formula \eqref{raw:conjugation} proves
$R^{-1}M(16,p)R=M(16,p)$; the same identity for $W$ follows from its Fricke
normalization.  Thus conjugation by every element of $G$ preserves this
Hecke set and permutes its left $\Gamma_0$-cosets.
These cosets have the usual representatives
\[
 \beta_\infty=\begin{pmatrix}p&0\\0&1\end{pmatrix},\qquad
 \beta_j=\begin{pmatrix}1&j\\0&p\end{pmatrix}
 \quad(0\leq j<p).
\]
For every $\beta\in M(16,p)$, the cusp $\beta(1/4)$ has denominator with
$2$-adic valuation exactly two after reduction to lowest terms: its
numerator is odd and its denominator is $4$ times an odd integer, and any
cancellation factor divides the odd integer $p$.  Hence it belongs to one
of the two denominator-$4$ classes.  Both have width one and belong to the
same $G$-orbit.  This is the cusp condition needed in the trace argument.

Let $P\in\mathbb Z[T]$.  Define
\[
 Q_P(\tau)=\frac{\Theta(\tau)P(x(p\tau))}{\Theta(p\tau)},\qquad
 G_P(q)=\frac{\Theta(q^p)P(x(q))}{\Theta(q)},
\]
and, using the determinant-normalized slash operator,
\[
 R_\beta(\tau)=p\,
       \frac{\Theta(\tau)}{(\Theta|_2\beta)(\tau)}P(x(\beta\tau)).
\]
The sum $\operatorname{Tr}_p(P)=\sum_{[\beta]}R_\beta$ is $G$-invariant.
Indeed left multiplication by $\Gamma_0$ does not change a summand;
right multiplication by $\Gamma_0$ permutes the summands; and for either
normalizer $g=R,W$, the factors $\Theta|_2g=\pm\Theta$ cancel to give
$R_\beta(g\tau)=R_{g^{-1}\beta g}(\tau)$.  Moreover
\begin{equation}\label{raw:trace}
 \operatorname{Tr}_p(P)=Q_P+p^3U_pG_P,
 \qquad U_p\left(\sum b_nq^n\right)=\sum b_{pn}q^n.
\end{equation}
For the $p$ finite representatives, the factor is $p^2$ before summing,
and the root-of-unity sum contributes the remaining factor $p$.

There are no poles in the upper half plane or at cusps outside the pole
orbit of $x$.  The latter assertion follows either directly from the
denominator classes, or from the adjugate argument: if $\beta c'$ were in
the pole orbit, then $p\beta^{-1}$ and normalizer conjugation would put
$c'$ in that orbit as well.
At $c=1/4$, precisely one Hecke branch multiplies the local order by $p$,
and the other $p$ branches divide it by $p$.  This follows from Hermite
normal form after a matrix sending $\infty$ to $c$.  Equality of the cusp
widths makes these the actual local scaling factors.  Since the orders of
$\Theta$ and $x$ are $2$ and $-2$, the unique possible largest pole has
order at most $2(p-1)+2p\deg P$.  Consequently
\begin{equation}\label{raw:tracepoly}
 \operatorname{Tr}_p(P)=\Pi_P(x),\qquad
 \Pi_P\in\mathbb Z[T],\qquad
 \deg\Pi_P\leq p\deg P+p-1.
\end{equation}
Polynomiality follows from the Hauptmodul and the single pole orbit.
Integrality follows by comparing the integral $q$-expansion in
\eqref{raw:trace}, using $x=q+O(q^2)$.

Let $q(X)\in X+X^2\mathbb Z[[X]]$ be the compositional inverse of $x(q)$,
and define
\[
 \sigma_p(X)=x(q(X)^p),\qquad
 \rho_p(X)=\mathcal F(X)/\mathcal F(\sigma_p(X)).
\]
The first series is congruent to $X^p$ modulo $p$.

\begin{theorem}\label{raw:formal}
For every odd prime $p$ and all $m,s\geq1$,
\[
 \rho_p(X)\mathcal F_{mp^{s-1}}(\sigma_p(X))
       -\mathcal F_{mp^s}(X)
       \in p^{3s}\mathbb Z_p[[X]].
\]
\end{theorem}

\begin{proof}
Fix $m$.  Set $K_j=mp^j$, and interpret $\mathcal F_{K_{-1}}$ as zero.
Apply \eqref{raw:trace}--\eqref{raw:tracepoly} to
$P=\mathcal F_{K_{s-1}}$.  Its trace polynomial has degree strictly less
than $K_s$, because $p(K_{s-1}-1)+p-1=K_s-1$.
Assuming the assertion at level $s-1$, its $q$-form gives
\[
 G_P(q)\equiv\mathcal F_{K_{s-2}}(x(q^p))
                    \pmod{p^{3(s-1)}}.
\]
Applying $p^3U_p$ gives $p^3\mathcal F_{K_{s-2}}(x(q))$ modulo $p^{3s}$.
For $s=1$ the same conclusion holds because both sides are zero modulo
$p^3$.  Thus $Q_P$ is congruent modulo $p^{3s}$ to an integral polynomial
in $x$ of degree less than $K_s$.
On the other hand
\[
 Q_P=\mathcal F(X)
       \frac{\mathcal F_{K_{s-1}}(\sigma_p(X))}
            {\mathcal F(\sigma_p(X))}
     \in\mathcal F(X)+X^{K_s}\mathbb Z[[X]].
\]
Comparing coefficients below $K_s$ identifies the polynomial modulo
$p^{3s}$ with $\mathcal F_{K_s}$.  Coefficients at or above $K_s$ vanish
modulo $p^{3s}$, as required.  This also starts the induction.
\end{proof}

\section{Integral specialization on the ordinary modular curve}\label{sec:ordinary}

We now restrict to $p\equiv1\pmod4$.  The use of the ordinary modular curve
in this section is essential: the second Hasse--Witt factor of the
cubic parameter vanishes at $64X^2=1$, so the original BV3 ring that
inverts that factor cannot be evaluated there.

Conjugating by $D_4$ identifies $X_0(16)$ analytically with the full-level
curve $X(4)$ (with $\pm1$ identified).  Since $p\equiv1\pmod4$, choose a
fourth root of unity in $\mathbb Z_p$; the corresponding component of
$X(4)$ has good reduction and a Tate cusp with local parameter
$q=e^{2\pi iz/4}$.  The functions and forms in \eqref{raw:theta} are
defined on this component.  They have coefficients in $\mathbb Z_p$.

For a precise integral chart, let $V$ be the smooth compactified
level-four component over $\mathbb Z_p$ with the two pole-cusp sections
removed, and let $A_0=\mathcal O(V)$. This is affine. The functions have
no vertical poles at odd primes: for $x$ use
\eqref{raw:lambda} and normality; for $\Theta$ use
\[
 \Theta^6=\frac{\Delta(4\tau)^5}
                  {\Delta(2\tau)^2\Delta(8\tau)^2}.
\]
The curves in this identity are related by prime-to-$p$ isogenies, and
their discriminants do not vanish on smooth elliptic fibers.
At retained cusps the leading constants are powers of two and roots
of unity. Consequently $\Theta$ trivializes the weight-two Hodge line
bundle on $V$ in characteristic zero and in characteristic $p$.

Let $\operatorname{Ha}$ denote the mod-$p$ Hasse invariant of weight
$p-1$. The quotient
$\operatorname{Ha}/\Theta^{(p-1)/2}$ is a function on $V_{\mathbb F_p}$.
Choose a lift $h_4\in A_0$ and put
\[
 A=A_0[1/h_4],\qquad \widehat A=\varprojlim_n A/p^nA.
\]
Its special fiber is the ordinary locus of $V$ and is geometrically
integral. At the Tate cusp $h_4$ is a unit, so $q$-expansions lie in
$\mathbb Z_p[[q]]$. The resulting formal open is stable under the
canonical quotient, since its reduction is Frobenius. A CM level
marking may require a finite unramified extension; all the following
arguments remain valid after that base change.

The following standard facts about the ordinary
modular curve are used in their integral form.

\begin{enumerate}
\item Quotienting by the canonical connected subgroup of order $p$ gives
an integral Frobenius lift $\Phi$ on the ordinary formal curve.  On the
Tate parameter it sends $q$ to $q^p$.  It carries prime-to-$p$ level
structures along the quotient and reduces to Frobenius modulo $p$.
\item If $\phi:E\to E'$ is the canonical isogeny, then
$p^{-1}\phi^*: \omega_{E'}\to\omega_E$ is an isomorphism on the ordinary
formal locus.  For the weight-two form $\Theta$, its ratio with the
image of $\Phi^*\Theta$ under this normalized map is an integral unit
$\rho\in\widehat A^\times$, with $q$-expansion
$\Theta(q)/\Theta(q^p)$.
\item The $q$-expansion homomorphism is injective modulo $p$ on $A/pA$.
Indeed a rational function on the integral special-fiber modular curve
whose expansion at a smooth cusp is zero is identically zero.
Consequently, for every $e\geq0$,
\begin{equation}\label{raw:saturation}
 \widehat A\cap p^e\mathbb Z_p[[q]]=p^e\widehat A.
\end{equation}
The higher powers follow by successive division by $p$.  This is the
integral $q$-expansion principle, not an assertion about convergence of
the individual hypergeometric period at a unit argument.
\end{enumerate}
These are the standard canonical-subgroup and $q$-expansion constructions
on the ordinary modular curve: see Katz, Proposition~2.7.1 for
divisibility detected by a $q$-expansion, Theorems~3.1 and~3.3 and
Section~3.4 for the ordinary canonical quotient, and Appendix~A2.1--A2.2
for the normalized differential and its Tate-curve expansion \cite{Katz}.  They can
also be made on a finer auxiliary prime-to-$p$ level and then descended.

The pole cusps and the ordinary locus are preserved under reduction of
$\Phi$, so $\Phi^*x$ is integral on this formal open.  Its $q$-expansion
is $x(q^p)$; therefore it realizes the same formal $\sigma_p$ as above.
All terms of Theorem~\ref{raw:formal} belong to $\widehat A$.  Applying
\eqref{raw:saturation} to their $q$-expansions yields
\begin{equation}\label{raw:analytic}
 \rho\,\mathcal F_{mp^{s-1}}(\Phi^*x)
       -\mathcal F_{mp^s}(x)\in p^{3s}\widehat A.
\end{equation}
Thus the full exponent $3s$ is retained under evaluation at every point
of this ordinary formal open.  No assertion about convergence of
$\mathcal F(1/8)$ is required.


\subsection{The ordinary CM point and its canonical branch}

At $z=i$, equivalently $\tau=i/4$, one has $\lambda(i)=1/2$ and
$x(i/4)=1/8$. The sign follows from the positive real theta constants.
The elliptic curve in the $z$ coordinate has complex multiplication by
$\mathbb Z[i]$; at $p\equiv1\pmod4$ it has good ordinary reduction.
A full level-four marking can be chosen over a finite unramified
extension with integer ring $\mathcal O$.

Write $p=a^2+b^2$, with $a$ odd and $b$ even, and choose the Gaussian
generator $\beta$ of the prime above $p$ which has positive valuation in
the chosen embedding. Its kernel has connected reduction, so the isogeny
given by $\beta$ is the canonical connected quotient at this CM lift.
Since $\beta\equiv1\pmod2$, it fixes the full level-two structure.
After the target is identified with the same CM elliptic curve, its
Legendre parameter is therefore still $\lambda=1/2$.

Applying $\Phi^*$ to $16x^2=\lambda(1-\lambda)$ now gives
\[
 (\Phi^*x)^2=1/64
\]
at the CM point. The integral Frobenius lift reduces to $x\mapsto x^p$,
so $\Phi^*x\equiv1/8\pmod p$. Since $p$ is odd, the two roots $1/8$
and $-1/8$ are distinct modulo $p$. Consequently
\begin{equation}\label{raw:fixed}
 (\Phi^*x)(i/4)=1/8.
\end{equation}
This proves fixedness on the required canonical branch.

For an explicit check on the finer level structure, let
$\alpha=a+bi$ be the unit generator and $\beta=a-bi$, so that
$\alpha\beta=p$. The matrix of multiplication by $\beta$ on level-four
homology, up to transpose and sign convention, reduces modulo $4$ to
$\pm I$ or
\[
 \pm R'\equiv\pm\begin{pmatrix}1&2\\2&1\end{pmatrix}.
\]
Indeed $a$ is odd, and $b$ is either $0$ or $2$ modulo $4$. Both matrices
preserve $x$ and have eigenvalue $+1$ on $\Theta$.
If $\omega$ is the differential after identifying the target with the
same CM curve, the normalized target differential is
$\omega'=(p/\beta)\omega=\alpha\omega$, since
$\phi^*\omega'=p\omega$. Weight two means that
$\Theta(E,\alpha\omega)=\alpha^{-2}\Theta(E,\omega)$. Hence
\begin{equation}\label{raw:scalar}
 \rho(i/4)=\alpha^2.
\end{equation}
The main conclusion below requires only that this value is a single
fixed unit $L_p$; its exact gamma normalization will follow from Guo's
published theorem.

All CM statements here concern the maximal-order curve at $z=i$.
The original lattice with period ratio $\tau=i/4$ has conductor four,
and an unrestricted norm-$p$ endomorphism assertion on that lattice
would be incorrect. The level-four model and the canonical quotient
give the fixedness statement actually used.


\section{The ordinary gamma multiplier and arbitrary-length recurrence}\label{sec:gamma}

Let $L_p$ be the value of $\rho$ at the ordinary CM point. The previous
section proves that $L_p$ is a fixed $p$-adic unit, independent of $m,s$,
and that for every $m,s\ge1$
\begin{equation}\label{eq:geometric}
 A_{m,s}\equiv L_pA_{m,s-1}\pmod{p^{3s}}.
\end{equation}
One may compute $L_p$ as the Gaussian unit-root square in
\eqref{raw:scalar}. For the gamma normalization, there is a shorter route
which uses only the published weaker recurrence.

\begin{lemma}\label{lem:units}
For $p\equiv1\pmod4$, each $H_s$ is a $p$-adic unit, and
$H_s\equiv H_1^s\pmod p$.
\end{lemma}
\begin{proof}
Put $d=(p-1)/2$. The coefficient
$c_k=\binom{2k}{k}^3/64^k$ is zero modulo $p$ if a base-$p$ digit of
$k$ exceeds $d$. Otherwise Lucas's congruence and Fermat's theorem give
$c_k\equiv\prod_j c_{k_j}\pmod p$, where $k_j$ are those digits.
Every digit string with $0\le k_j\le d$ represents an integer at most
$(p^s-1)/2$. Hence the half-length sum factors as $H_1^s$ modulo $p$.
Guo's published first-level case gives
$H_1\equiv-\Gamma_p(1/4)^4\pmod p$, which is a unit.
\end{proof}

\begin{lemma}\label{lem:scalar}
The constant $L_p$ in \eqref{eq:geometric} equals
$-\Gamma_p(1/4)^4$ exactly in $\mathbb Z_p$.
\end{lemma}
\begin{proof}
Write $g_p=-\Gamma_p(1/4)^4$. Guo's theorem
\cite[arXiv equation~(1.9)]{Guo} states
\[
 H_s\equiv g_pH_{s-1}\pmod{p^{s+1}}\qquad(s\ge1).
\]
The case $m=1$ of \eqref{eq:geometric} gives the same recurrence with
$L_p$ and modulus $p^{3s}$. Subtracting and using Lemma~\ref{lem:units}
therefore gives $L_p-g_p\in p^{s+1}\mathbb Z_p$ for every $s$.
The intersection of these ideals is zero. This proves the exact equality.
If the CM marking was evaluated over an unramified extension with
integer ring $\mathcal O$, run this argument first in $\mathcal O$.
It proves the same equality and hence descent to $\mathbb Z_p$.
\end{proof}

\begin{proof}[Proof of the ordinary assertion in Theorem~\ref{thm:main}]
Substitute Lemma~\ref{lem:scalar} into \eqref{eq:geometric}. This gives
the assertion with modulus $p^{3s}$. Write $m=p^v m_0$, with $p\nmid m_0$.
Then $A_{m,s}=A_{m_0,s+v}$ and
$A_{m,s-1}=A_{m_0,s+v-1}$. Applying the same assertion at level $s+v$
gives the displayed stronger modulus $p^{3s+3v_p(m)}$.
For $m=1$ the upper summation index is exactly $(p^s-1)/2$; for $m=2$
it is $p^s-1$. Thus both stated truncation conventions follow.
\end{proof}


\section{Integral Legendre jets and change of Frobenius lift}\label{sec:jets}

Throughout this section, $p$ may be any odd prime. The local ring
notation is reset on the punctured Legendre base; the restriction
$p\equiv1\pmod4$ in Section~\ref{sec:ordinary} is not imposed here.
Write
\[
 D=\lambda(1-\lambda),\qquad h=1-2\lambda,\qquad
 A=\widehat{\mathbb Z_p[\lambda,D^{-1}]},
\]
where completion is $p$-adic.  Put
\[
 N_s=(p^s-1)/2,\qquad
 P_s(\lambda)=\sum_{k=0}^{N_s}\binom{2k}{k}^3(D/16)^k,
 \qquad Y_s=(P_s,P_s',P_s'')^{\mathsf T}.
\]
The desired half-length sum is
\[
 H_s=P_s(1/2)=\sum_{k=0}^{(p^s-1)/2}\frac{(1/2)_k^3}{k!^3}.
\]
The quadratic pullback is essential: the differential equation in the
raw cubic parameter is singular at $X=1/8$, whereas the connection used
below is regular at $\lambda=1/2$.

The formal input is Theorem~\ref{raw:formal}, proved above by the raw Hecke trace.  In the present notation, on the ordinary locus it gives
\begin{equation}\label{ss:ordinaryscalar}
 P_s(\lambda)-\rho(\lambda)P_{s-1}(\sigma(\lambda))
       \in p^{3s}A^{\mathrm{ord}},
\end{equation}
where $\sigma$ is the ordinary canonical quotient lift and
$\rho=(L/L^\sigma)^2$, with
$L={}_2F_1(1/2,1/2;1;\lambda)$.  This is an identity of integral analytic
quotients; the individual series $L$ need not converge at unit arguments.
Here
\[
 A^{\mathrm{ord}}=
 \widehat{\mathbb Z_p[\lambda,D^{-1},\mathscr H^{-1}]}
\]
and $\mathscr H$ is any integral lift of the nonzero Legendre Hasse
polynomial; both hats denote $p$-adic completion.  Equivalently it is
the $p$-adic completion of the localization of $A$ at $\mathscr H$.
To obtain \eqref{ss:ordinaryscalar}, substitute
$X^2=D/16$ in the even raw truncations and use the ordinary
$q$-expansion principle.  The corresponding ordinary Frobenius lifts
commute with this modular pullback.  All terms are integral on the
ordinary locus, and divisibility of their cusp expansions gives the
stated divisibility in $A^{\mathrm{ord}}$; the integral ordinary
$q$-expansion principle is Katz's Proposition~2.7.1
\cite{Katz}, and the canonical quotient is constructed in his
Theorems~3.1 and~3.3 and Section~3.4.  The raw Tate parameter satisfies
$q_E=q_{\rm raw}^{\,4}$, so its replacement by $q_{\rm raw}^p$ is exactly
$q_E\mapsto q_E^p$ on the Legendre Tate curve.  The fact that the raw
truncations are even makes the substitution $X^2=D/16$ integral without
choosing a square root of $\lambda$.

\subsection{An exact differential-equation residual}

Set
\[
 P=\frac{h}{D},\qquad Q=-\frac1{4D}.
\]
The Legendre equation is $y''+Py'+Qy=0$.  Its symmetric-square equation
has monic operator
\[
 \mathcal L=\partial^3+3P\partial^2+
    (P'+4Q+2P^2)\partial+(2Q'+4PQ).
\]
Equivalently,
\begin{equation}\label{ss:operator}
 D^2\mathcal L=D^2\partial^3+3Dh\partial^2+
                         (1-7D)\partial-h/2.
\end{equation}

\begin{lemma}\label{ss:residual}
For every $s\geq0$,
\[
 \mathcal L P_s
  =-\frac{p^{3s}}2\,
       \frac{\binom{2N_s}{N_s}^{3}}{16^{N_s}}
       h D^{N_s-2}\in p^{3s}A.
\]
In particular the jet vector satisfies
$Y_s'=CY_s+r_se_3$, where $C$ is the integral companion matrix of
$\mathcal L$, $r_s\in p^{3s}A$, and $e_3=(0,0,1)^{\mathsf T}$.
\end{lemma}

\begin{proof}
Direct differentiation gives
\[
 (D^2\mathcal L)(D^k)
   =h\left(k^3D^{k-1}-\frac{(2k+1)^3}{2}D^k\right).
\]
For $a_k=\binom{2k}{k}^3/16^k$ one has
$k^3a_k=(2k-1)^3a_{k-1}/2$.  All interior terms therefore cancel, and
the last term is $-(2N_s+1)^3a_{N_s}hD^{N_s}/2$.
Since $2N_s+1=p^s$, division by the unit $D^2$ in $A$ proves the claim.
\end{proof}

\subsection{The integral rank-three Frobenius matrix}

We use the constant quadratic twist
\[
 E^-_\lambda:\quad y^2=-x(x-1)(x-\lambda).
\]
At $\lambda=0$ its nodal tangents are $\pm1$, so its Tate degeneration
is split.  This makes the following weight-one normalization unambiguous.
For the usual Legendre curve the weight-one matrix differs by the
constant sign $(-1/p)$; its symmetric square is identical.  This sign
and the usual Legendre Frobenius normalization are recorded explicitly
in Shirai's introduction, equations~(1), (3), and~(4) \cite{Shirai}.

Let $\omega=dx/y$ and $\eta=\nabla_{\partial_\lambda}\omega$.
These form an integral basis of $H^1_{\mathrm{dR}}(E^-_\lambda/A)$.
Indeed in the usual basis $(\omega,x\omega)$ one has
\begin{equation}\label{ss:basis}
 \eta=\frac{\omega}{2(1-\lambda)}
       -\frac{x\omega}{2\lambda(1-\lambda)},
\end{equation}
whose determinant is a unit for every odd $p$.  The relation follows by
reducing the exact differential $d(x(x-1)/y)$.  The Gauss--Manin
connection is
\[
 \nabla\binom{\omega}{\eta}
 =\begin{pmatrix}0&1\\-Q&-P\end{pmatrix}
       \binom{\omega}{\eta}\,d\lambda.
\]

Let $\phi:A\to A$ be any lift of Frobenius.  The integral crystalline
Frobenius of this smooth proper elliptic family gives a matrix $M_\phi$
with entries in $A$.  Our convention is a row convention:
\[
 F_\phi(\phi^*(\omega,\eta)^{\mathsf T})
       =M_\phi(\omega,\eta)^{\mathsf T}.
\]
Its determinant is $p$ times a unit, by the cup product and the action
of Frobenius on $H^2$.  Consequently
\begin{equation}\label{ss:V}
 V_\phi=pM_\phi^{-1}\in M_2(A).
\end{equation}
This uses the adjugate of a rank-two matrix and does not invert the
Hasse invariant.  For the underlying integral comparison, one can use
the Dieudonne crystal of the elliptic scheme's $p$-divisible group:
Berthelot--Messing \cite[\S2, Theorem~2, p.~23]{BM} gives its local
freeness, the construction on pp.~24--25 gives integral $F,V$ with
$FV=VF=p$, and \cite[\S3, equation~(3.1), p.~27]{BM} identifies it with
relative $H^1_{\rm cris}$, compatibly with the Hodge filtration.
Its evaluation on the smooth lift is integral de Rham cohomology.
For the completed comparison and the canonical integral Frobenius map,
see also \cite[Remark~60.24.14, Tag~07MZ; Theorem~60.26.4,
Tag~07N5]{Stacks}; the latter map is constructed before inverting $p$.
The first Hodge line is carried into $pH^1$ by Frobenius: modulo $p$
this follows from the pullback of regular differentials under relative
Frobenius being zero, together with the integral comparison and its
functoriality.  In the row convention this says that the first row
of $M_\phi$ is divisible by $p$ at a Frobenius-fixed fiber.

The symmetric-square jet basis is
\[
 \omega^2,\qquad2\omega\eta,\qquad
 2\eta^2-2P\omega\eta-2Q\omega^2.
\]
Its change-of-basis matrix from $(\omega^2,\omega\eta,\eta^2)$ is
\begin{equation}\label{ss:S}
 S=\begin{pmatrix}1&0&0\\0&2&0\\-2Q&-2P&2\end{pmatrix},
 \qquad\det S=4.
\end{equation}
Define the integral period-side Frobenius matrix
\begin{equation}\label{ss:Bphi}
 B_\phi=S\operatorname{Sym}^2(V_\phi)(S^\phi)^{-1}\in M_3(A).
\end{equation}
The word ``period-side'' records the inverse in \eqref{ss:V}; confusing
this matrix with the cohomology-row matrix would reverse the required
Hodge divisibility.

On the ordinary canonical lift $\sigma$, write $a=L/L^\sigma$.
Here the exact differential normalization deserves to be specified.
The connected subgroup quotient reduces to relative Frobenius;
transporting the ordered $2$-torsion gives the parameter $\sigma(\lambda)$.
Its model-scaling unit relative to $E^-_{\sigma(\lambda)}$ is congruent
to one modulo $p$ and therefore has a unique square root congruent to
one for odd $p$.  This identifies the quotient pullback with the
crystalline $F_\sigma$, without an unspecified quadratic twist.
One may make this identification on a finite etale auxiliary level
and descend; the matrices in our displayed basis already lie over $A$.

On the split Tate chart, the differential satisfies
$\omega=L(\lambda)\omega_{\rm Tate}$, with constant term one.  This
identity follows either from the normalized Tate period or from its
Legendre differential equation and its initial value at $\lambda=0$.
The canonical isogeny sends $z\mapsto z^p$ and hence pulls back
$\omega_{\rm Tate}$ to $p\omega_{\rm Tate}$
\cite[Appendix~A2, especially (A2.2.1) and (A2.2.7), pp.~175--177]{Katz}.
Consequently its differential is exactly
\[
 \omega^\sigma\longmapsto c\omega,
 \qquad c=p\frac{L^\sigma}{L}=\frac p a.
\]
The identity on the whole ordinary ring follows from the integral
$q$-expansion principle; thus no parameter-dependent unit is being
suppressed.  Horizontality now gives
\[
 M_\sigma=
 \begin{pmatrix}c&0\\c'/\sigma'&c/\sigma'\end{pmatrix},
 \qquad
 pM_\sigma^{-1}=
 \begin{pmatrix}a&0\\a'&a\sigma'\end{pmatrix}.
\]
These equalities can be read in the fraction field; the resulting
matrices are the integral ones above.  Taking the symmetric square in
the jet basis yields
\begin{equation}\label{ss:Bsigma}
 B_\sigma=
 \begin{pmatrix}
 \rho&0&0\\
 \rho'&\rho\sigma'&0\\
 \rho''&2\rho'\sigma'+\rho\sigma''&\rho(\sigma')^2
 \end{pmatrix},\qquad \rho=a^2.
\end{equation}
In particular its columns have respective valuations at least $0,1,2$,
because $\sigma'\equiv\sigma''\equiv0\pmod p$.
Differentiating \eqref{ss:ordinaryscalar} twice therefore gives
\begin{equation}\label{ss:ordinaryjet}
 Y_s-B_\sigma Y_{s-1}^\sigma\in(p^{3s}A^{\mathrm{ord}})^3.
\end{equation}

\subsection{Changing the lift without losing three digits}

Recall the definition of $\epsilon_p$ from Section~\ref{sec:statement}.

\begin{lemma}[Taylor transport with a third-order residual]
\label{ss:taylor}
Let $R$ be one of the above $p$-adically complete differential rings, with
$\partial R\subseteq R$, and let $u,v:R\to R$ be substitutions congruent
modulo $p$.  Suppose $Y=(f,f',f'')^{\mathsf T}$ and
$Y'=CY+re_3$, with $r\in p^bR$ and the integral companion matrix $C$.
If $T(u,v)$ is horizontal Taylor transport from $v$ to $u$, then
\[
 Y^u-T(u,v)Y^v\in
 p^{b+3-\epsilon_p}R\ \oplus\
 p^{b+2}R\ \oplus\ p^{b+1}R.
\]
\end{lemma}

\begin{proof}
Put $\delta=u(\lambda)-v(\lambda)\in pR$.
The matrices of iterated covariant differentiation have integral
entries.  Differentiating the inhomogeneous companion system repeatedly
shows that its residual first enters the three jet coordinates at
derivative orders $3,2,1$, respectively.  Subsequent coefficients are
integral combinations of the coefficients of $C$, their derivatives,
and derivatives of $r$; hence they remain in $p^bR$.

Subtract the horizontal Taylor series from the ordinary Taylor series
for $Y$.  The resulting three sums start at orders $3,2,1$ and their
$n$th terms are divisible by $p^b\delta^n/n!$.  They converge since
$v_p(\delta^n/n!)\geq n-v_p(n!)\to\infty$ for odd $p$.
For the three starting orders, the respective lower bounds are
\[
 \min_{n\geq3}(n-v_p(n!))=3-\epsilon_p,\quad
 \min_{n\geq2}(n-v_p(n!))=2,\quad
 \min_{n\geq1}(n-v_p(n!))=1.
\]
This proves the assertion and identifies the unique extra loss at $p=3$.
\end{proof}

\begin{proposition}\label{ss:regularjet}
For any integral Frobenius lift $\phi$ of the punctured Legendre base and
all $s\geq1$,
\[
 Y_s-B_\phi Y_{s-1}^\phi\in(p^{3s-\epsilon_p}A)^3.
\]
\end{proposition}

\begin{proof}
First work on the ordinary localization.  The crystalline
change-of-lift isomorphism is horizontal Taylor transport; see
\cite[\S2.1, pp.~19--20]{BM} and
\cite[Proposition~60.17.4, Tag~07JH]{Stacks}, where the divided-power
Taylor formula and its cocycle identity are explicit.
In the rank-two row convention, if $T_2(\sigma,\phi)$ transports
solution coordinates from $\phi$ to $\sigma$, the cohomology matrices
satisfy
\[
 M_\phi=T_2(\sigma,\phi)^{-1}M_\sigma,
 \qquad V_\phi=V_\sigma T_2(\sigma,\phi).
\]
Taking symmetric squares and changing to the jet bases gives precisely
$B_\phi=B_\sigma T(\sigma,\phi)$.
Apply Lemma~\ref{ss:taylor} to $Y_{s-1}$, using
$b=3(s-1)$ from Lemma~\ref{ss:residual}.  Multiplication by
$B_\sigma$ adds at least $0,1,2$ to its three columns.  Therefore
\[
 B_\sigma\bigl(Y_{s-1}^\sigma-
                    T(\sigma,\phi)Y_{s-1}^\phi\bigr)
                \in(p^{3s-\epsilon_p}A^{\mathrm{ord}})^3.
\]
Combine this with \eqref{ss:ordinaryjet}.

Both $Y_s$ and $B_\phi Y_{s-1}^\phi$ already belong to $A^3$, since
\eqref{ss:Bphi} is integral on the whole smooth Legendre base.
The map $A/pA\to A^{\mathrm{ord}}/pA^{\mathrm{ord}}$ is injective:
the latter is localization at the nonzero Hasse polynomial in the
integral domain $\mathbb F_p[\lambda,D^{-1}]$.
Successive division by $p$ then gives
$A\cap p^eA^{\mathrm{ord}}=p^eA$ for every $e$.
This descends the congruence without any additional loss.
\end{proof}

\section{Supersingular CM specialization}\label{sec:sscm}

Assume now $p\equiv3\pmod4$ and choose
\[
 \phi(\lambda)=\frac12+\left(\lambda-\frac12\right)^p.
\]
It is an integral Frobenius lift, preserves invertibility of $D$
modulo $p$, and fixes $\lambda_0=1/2$.
The elliptic curve $E^-_{1/2}$ has good supersingular reduction and
Frobenius trace zero.  For example, it is a quadratic twist of
$y^2=x^3-x$; at $p\equiv3\pmod4$ the trace of that curve is zero by
pairing $x$ with $-x$ in its quadratic-character sum.
Let $M=M_\phi(1/2)$.  Its characteristic polynomial is $T^2+p$, so
\[
 M^2=-pI_2,\qquad V=pM^{-1}=-M.
\]
Frobenius--Hodge divisibility says that the first row of $M$ is in
$p\mathbb Z_p^2$.  Consequently, for $B=B_\phi(1/2)$,
\begin{equation}\label{ss:CMmatrix}
 B^2=p^2I_3,\qquad (1,0,0)B\in p^2\mathbb Z_p^3.
\end{equation}
The first assertion follows by taking a symmetric square.  For the
second, the first symmetric-square basis vector is $\omega^2$, and the
first row of $S$ in \eqref{ss:S} is $(1,0,0)$; multiplying by the integral
inverse of $S$ preserves divisibility by $p^2$.

\begin{theorem}\label{ss:main}
For every prime $p\equiv3\pmod4$ and every integer $s\geq2$,
\[
 H_s-p^2H_{s-2}\in p^{3s-1-\epsilon_p}\mathbb Z_p.
\]
Thus the exponent is $3s-1$ for $p\geq7$, and $3s-2$ for $p=3$.
\end{theorem}

\begin{proof}
Evaluate Proposition~\ref{ss:regularjet} at $\lambda_0$ and write
\[
 Y_s(\lambda_0)=B Y_{s-1}(\lambda_0)+E_s,
 \qquad E_s\in(p^{3s-\epsilon_p}\mathbb Z_p)^3.
\]
Applying this at two consecutive levels gives
\[
 Y_s(\lambda_0)-p^2Y_{s-2}(\lambda_0)=E_s+BE_{s-1}.
\]
By \eqref{ss:CMmatrix}, the first coordinate of the second term is
divisible by $p^{2+3(s-1)-\epsilon_p}=p^{3s-1-\epsilon_p}$; the first
term has one more factor of $p$.  The first coordinate of $Y_s(\lambda_0)$
is $H_s$, proving the asserted two-step recurrence.
\end{proof}

\subsection{A coupled auxiliary-sum formulation}

For $s\geq0$ put
\[
 W_s=\sum_{k=0}^{(p^s-1)/2}(4k+1)\frac{(1/2)_k^3}{k!^3}
     =H_s-\frac12P_s''(1/2).
\]
The second equality follows by differentiating $D=\lambda(1-\lambda)$
twice at $\lambda=1/2$: $D'=0$, $D''=-2$, and $D=1/4$.

\begin{corollary}\label{ss:coupled}
For each $p\equiv3\pmod4$ there is a unit
$\mu_p\in\mathbb Z_p^\times$, the square of an integral crystalline
matrix coefficient, such that for every $s\geq1$,
\[
 \begin{split}
 H_s&\equiv-\frac{p^2}{\mu_p}W_{s-1}
                    \pmod{p^{3s-\epsilon_p}},\\
 W_s&\equiv-\mu_p H_{s-1}
                    \pmod{p^{3s-\epsilon_p}}.
 \end{split}
\]
No special-function identification of $\mu_p$ is needed or asserted.
\end{corollary}

\begin{proof}
At $\lambda=1/2$, equation~\eqref{ss:basis} gives
$\eta=(1-2x)\omega$.  After adjoining $i$, the CM automorphism
$(x,y)\mapsto(1-x,iy)$ acts on $(\omega,\eta)$ as
$\operatorname{diag}(i,-i)$.  Since $p\equiv3\pmod4$, crystalline
Frobenius conjugates $i$ to $-i$.  Its matrix, and hence $V=-M$, is
off diagonal:
\[
 V=\begin{pmatrix}0&b\\c&0\end{pmatrix},\qquad bc=-p.
\]
The Hodge-row divisibility and integrality give $v_p(b)=1$ and
$v_p(c)=0$.  Set $\mu_p=c^2$; then $b^2=p^2/\mu_p$.

At this point $P=0$ and $Q=-1$, so in symmetric-square solution
coordinates the first and third entries of the vector corresponding
to $Y_s$ are
\[
 H_s,\qquad \tfrac12P_s''(1/2)-H_s=-W_s.
\]
The middle entry is zero because $P_s(\lambda)=P_s(1-\lambda)$.
The symmetric square of the displayed off-diagonal matrix interchanges
the first and third coordinates with factors $b^2$ and $c^2$.
Proposition~\ref{ss:regularjet}, with the integral basis change
\eqref{ss:S}, therefore gives the two congruences with the same modulus.
\end{proof}

In particular these congruences prove the exact valuations
\[
 v_p(H_s)=2\lceil s/2\rceil,\qquad
 v_p(W_s)=2\lfloor s/2\rfloor.
\]
For $p\geq7$ this follows immediately by induction from
$H_0=W_0=1$, because each error term has strictly higher valuation
than the displayed leading term.  For $p=3$ the same induction starts
from the explicit values $H_1=9/8$ and $W_1=13/8$; at all later levels
the error valuation is again strictly higher.
In particular $H_1\in p^2\mathbb Z_p$. Together with
Theorem~\ref{ss:main}, this proves the supersingular assertion
in Theorem~\ref{thm:main}, including its first level.

\paragraph{What supplies the exponent.}
The scalar ordinary theorem alone cannot be evaluated at this point.
The proof uses its differentiated rank-three form, the exact
$p^{3s}$ differential-equation residual, and weighted Taylor transport
to an integral crystalline matrix that is defined on the supersingular
fiber.  The loss of one digit comes from combining the previous-level
error $p^{3s-3}$ with the Hodge-row factor $p^2$.  The additional loss at
$p=3$ comes from $3!$ in the order-three Taylor term.  Frobenius slopes
or the identity $B^2=p^2I$ by themselves do not supply either estimate.

\section{Precision, scope, and exact computational checks}\label{sec:checks}

Theorem~\ref{thm:main} addresses the half-length sums at every stated
level for all odd primes, with exponent $3s-2$ at $p=3$.
For ordinary primes it also gives arbitrary raw lengths, including
the full cutoff. The supersingular proof uses
$2N_s+1=p^s$ in its exact differential residual; an extension to another
cutoff would require a separate residual estimate.

The two losses in the supersingular exponent have different sources.
The one-step jet congruence has precision $3s-\epsilon_p$.
In the two-step composition, the preceding error has precision
$3(s-1)-\epsilon_p$, and the first-row Hodge factor restores two digits.
This gives $3s-1-\epsilon_p$. The extra loss at $p=3$ occurs in
$\delta^3/3!$ in Taylor transport. Both estimates are part of the
proof, independently of the numerical patterns below.

\subsection{Ordinary half- and full-cutoff checks}

The ordinary verification program checks the multiplier as the
Gaussian unit root of
\[
 X^2-2(a^2-b^2)X+p^2,
 \qquad p=a^2+b^2,\quad a\text{ odd}.
\]
It checks $23$ prime/level pairs, in both half and full cutoffs,
at precision $p^{3s+3}$. Every tested half- and full-cutoff defect
has valuation exactly $3s$.
\begin{center}
\begin{tabular}{@{}lll@{}}\toprule
Prime & Levels $s$ & Observed half-sum defect valuations\\\midrule
$5$ & $1,\ldots,6$ & $3,6,9,12,15,18$\\
$13,17$ & $1,\ldots,4$ & $3,6,9,12$\\
$29,37,41$ & $1,2,3$ & $3,6,9$\\\bottomrule
\end{tabular}
\end{center}
At $p=5$, the upper-half tail at level two has valuation four, below
the target six, whereas its recurrence defect has valuation seven
in the exact calculation. The raw-cutoff argument handles both sums
without requiring that tail itself to vanish at the target precision.

\subsection{Supersingular half-cutoff checks}

The supersingular program works at precision $p^{3s_{\max}+5}$ and
removes $p$-power factors before every division. The table lists
$v_p(H_s-p^2H_{s-2})$ in increasing order of $s$. Each displayed
valuation is exact at the working precision.
\begin{center}
\begin{tabular}{@{}lll@{}}\toprule
Prime & Levels $s$ & Observed two-step defect valuations\\\midrule
$3$ & $2,\ldots,10$ & $4,7,10,13,16,19,22,25,28$\\
$7$ & $2,\ldots,6$ & $5,8,11,14,17$\\
$11$ & $2,\ldots,5$ & $5,8,11,14$\\
$19,23$ & $2,3,4$ & $5,8,11$\\
$31,43,47$ & $2,3$ & $5,8$\\\bottomrule
\end{tabular}
\end{center}
These are $30$ prime/level recurrence checks. They establish sharpness
in the tested cases, not equality of the defect valuation for every
prime and level.

\subsection{Formal checks and reproducibility}

A third program checks the raw formal identity directly from the eta
products. Its $13$ parameter sets include primes $3,5,7,11$, several
values of $m$, and levels through $s=4$. Computations use exact modular
arithmetic, with an explicit bound preventing integer overflow in
the NumPy convolutions. Each parameter set checks only the finite
coefficient window specified in the program.

All three programs are reproduced in Appendix~\ref{app:programs}.
The two numerical recurrence programs use only Python's standard
library; the eta-product program also requires NumPy. The combined
evidence consists of $23$ ordinary prime/level pairs (each checked in
two cutoffs), $30$ supersingular recurrence pairs, and $13$ finite
formal parameter sets. These categories are distinct. The general
claims are supplied by the preceding arguments, rather than by an
extrapolation from the computations.


\appendix
\section{An independent finite identity and valuation bound}\label{app:finite}

Fix an odd prime $p\equiv3\pmod4$.  Write
\[
 b_j=\frac{(1/2)_j}{j!}=\frac{\binom{2j}{j}}{4^j},\qquad
 q_s=p^s,\qquad n_s=(q_s-1)/2,\qquad
 H_s=\sum_{k=0}^{n_s}b_k^3.
\]
For an odd positive integer $t\equiv1\pmod4$, put
\[
 C(t)=b_{(t-1)/4}^2,
 \qquad U_s=\frac{\binom{p^s-1}{(p^s-1)/2}}{4^{p^s-1}}.
\]
This appendix supplies an independent finite identity and valuation
bound. It is not needed for the supersingular argument in
Sections~\ref{sec:jets}--\ref{sec:sscm}.

\subsection{Exact interpolation identity}

\begin{proposition}\label{prop:finite-lagrange}
For every integer $n\ge0$, with $q=2n+1$,
\begin{align}
 \sum_{k=0}^n b_k^3
 &=q^2b_n^2\sum_{l=0}^{\lfloor n/2\rfloor}
 \frac{b_l^2(n!)^2}{(4l+1)(n-2l)!(n+2l+1)!} \label{eq:lagrange-first}\\
 &=q\frac{\binom{2n}{n}}{16^n}
 \sum_{l=0}^{\lfloor n/2\rfloor}
 \binom{q}{n-2l}\frac{b_l^2}{4l+1}. \label{eq:lagrange-second}
\end{align}
\end{proposition}
\begin{proof}
Consider the polynomial in $T^2$ of degree $n$
\[
 \mathcal A_n(T)=\sum_{k=0}^n b_k^3
       \prod_{j=0}^{k-1}\left(1-\frac{T^2}{(2j+1)^2}\right).
\]
Let $\mathsf P_j$ denote the classical Legendre polynomial of degree $j$.
For $0\le j\le n$, the terminating Clausen identity gives
\[
 \mathcal A_n(2j+1)
 ={}_3F_2\!\left(\begin{matrix}-j,j+1,1/2\\1,1\end{matrix};1\right)
 =\mathsf P_j(0)^2.
\]
This value is zero for odd $j$, and is $b_{j/2}^2$ for even $j$.
Lagrange interpolation at the nodes $(2j+1)^2$ gives the coefficient
\[
 L_j(0)=\prod_{\substack{0\le i\le n\\i\ne j}}
 \frac{-(2i+1)^2}{(2j+1)^2-(2i+1)^2}
 =\frac{(-1)^j((2n+1)!!)^2}
 {4^n(2j+1)(n-j)!(n+j+1)!}.
\]
Only $j=2l$ contributes.  Since
$((2n+1)!!)^2/4^n=q^2b_n^2(n!)^2$, this proves
\eqref{eq:lagrange-first}.  Using
$q^2b_n^2(n!)^2/q!=q\binom{2n}{n}/16^n$ proves
\eqref{eq:lagrange-second}.
\end{proof}

In particular, define the finite index set
\[
 \mathcal T_s=\{t\in\mathbb Z:0<t\le p^s,\ t\equiv1\pmod4\}.
\]
Then the exact identity becomes
\begin{equation}\label{eq:exact-Hs-kernel}
 H_s=\sum_{t\in\mathcal T_s}\mathcal K_s(t),\qquad
 \mathcal K_s(t)=p^sU_s\binom{p^s}{(p^s-t)/2}\frac{C(t)}t.
\end{equation}
Every quantity in this equation is rational; no convergence or analytic
continuation is used.

\subsection{Carries and unconditional divisibility}

\begin{lemma}\label{lem:kernel-valuation}
If $t=4l+1\in\mathcal T_s$ and $a=v_p(t)$, then
\[
 v_p(\mathcal K_s(t))=2s+2v_p\binom{2l}{l}-2a,
 \qquad
 v_p\binom{2l}{l}\ge\lceil a/2\rceil.
\]
Consequently,
\[
 v_p(\mathcal K_s(t))\ge2s-2\lfloor a/2\rfloor,
 \qquad
 v_p(H_s)\ge2\lceil s/2\rceil.
\]
For even $s$, equality holds: $v_p(H_s)=s$.
\end{lemma}
\begin{proof}
$U_s$ is a $p$-adic unit.  If $t<p^s$ and
$k=(p^s-t)/2$, then $v_p(k)=a$ and
$v_p\binom{p^s}{k}=s-a$.  The latter follows from
$\binom{p^s}{k}=(p^s/k)\binom{p^s-1}{k-1}$ and Lucas's
theorem.  For $t=p^s$, the same formula gives $s-a=0$.
Equation \eqref{eq:exact-Hs-kernel} gives the exact valuation asserted.

For the carry bound, write $p=4h+3$.  The base-$p$ digits forced by
$4l+1\equiv0\pmod{p^a}$ alternate, from the least significant end,
between $3h+2$ and $h$.  If $a$ is odd, the last forced digit is
$3h+2$; if $a$ is even, it is $h$.  In adding $l+l$, each
$3h+2$ digit produces a carry, and each following $h$ digit,
with its incoming carry, produces none.  Thus at least
$\lceil a/2\rceil$ carries occur.  Kummer's theorem proves the bound.

For odd $s$, $p^s\equiv3\pmod4$, so $p^s\notin\mathcal T_s$ and
$a\le s-1$.  This gives valuation at least $s+1$.
For even $s$, the only term with $a=s$ is $t=p^s$, and its carry
count is exactly $s/2$.  Its valuation is $s$, whereas every other
term has valuation at least $s+2$.  There can be no cancellation of
this unique least-valued term.
\end{proof}



\section{Reproducible verification programs}\label{app:programs}
Each listing is a separate executable source file. Save it under the
specified name and run it with Python 3; it writes a JSON record next
to its source. Listings B.1 and B.3 use only the standard library.
Listing B.2 additionally requires NumPy. Their finite scope is described
in Section~\ref{sec:checks}.

\subsection{Ordinary half- and full-cutoff recurrences}
Save as \path{check_recurrences.py}.
\begin{lstlisting}
"""Finite exact checks; the all-level theorem is proved in the paper."""
from math import isqrt, comb
from pathlib import Path
import json

def strip(n, p):
    v = 0
    while n % p == 0:
        n //= p
        v += 1
    return n, v

def valuation(n, p, precision):
    n %= p**precision
    if n == 0:
        return {'value': precision, 'exact': False}
    return {'value': strip(n, p)[1], 'exact': True}

def terms(p, length, precision):
    """Terms ((1/2)_k/k!)^3; all divisions are on p-units."""
    modulus = p**precision
    unit, v = 1, 0
    for k in range(length):
        assert v >= 0
        yield (0 if 3*v >= precision else
               pow(unit, 3, modulus)*pow(p, 3*v, modulus) % modulus)
        a, va = strip(2*k+1, p)
        b, vb = strip(2*k+2, p)
        unit = unit*a*pow(b, -1, modulus) % modulus
        v += va-vb

def gaussian_trace(p):
    for a in range(1, isqrt(p)+1, 2):
        b = isqrt(p-a*a)
        if a*a+b*b == p:
            return a, b, 2*(a*a-b*b)
    raise ValueError(p)

def unit_root(p, trace, precision):
    modulus = p**precision
    u = trace % p
    assert u != 0
    for _ in range(precision):
        u = (u-(u*u-trace*u+p*p)*
             pow(2*u-trace, -1, modulus)) % modulus
    assert (u*u-trace*u+p*p) % modulus == 0
    return u

def check():
    records = []
    for p in [5, 13, 17, 29, 37, 41]:
        a, b, trace = gaussian_trace(p)
        for s in range(1, 7):
            if p**s > 100000:
                break
            precision = 3*s+3
            modulus = p**precision
            u = unit_root(p, trace, precision)
            current = list(terms(p, p**s, precision))
            previous = list(terms(p, p**(s-1), precision))
            half = sum(current[:(p**s+1)//2]) % modulus
            prevhalf = sum(previous[:(p**(s-1)+1)//2]) % modulus
            full = sum(current) % modulus
            prevfull = sum(previous) % modulus
            if s == 1:
                direct = sum(pow(comb(2*k, k), 3, modulus)*
                             pow(64, -k, modulus) for k in range(p)) % modulus
                assert full == direct
            rec = {'p': p, 's': s, 'a': a, 'b': b,
                   'precision': precision, 'modulus': modulus, 'unit_root': u,
                   'half': half, 'previous_half': prevhalf,
                   'full': full, 'previous_full': prevfull,
                   'half_defect': valuation(half-u*prevhalf, p, precision),
                   'full_defect': valuation(full-u*prevfull, p, precision),
                   'tail': valuation(full-half, p, precision),
                   'tail_defect': valuation(full-half-u*(prevfull-prevhalf),
                                             p, precision)}
            assert rec['half_defect'] == {'value': 3*s, 'exact': True}
            assert rec['full_defect'] == {'value': 3*s, 'exact': True}
            records.append(rec)
    return {'scope': 'Finite checks, not a proof of untested cases.',
            'prime_level_pairs': len(records), 'all_passed': True,
            'all_half_and_full_defects_exactly_3s': True, 'records': records}

if __name__ == '__main__':
    result = check()
    Path(__file__).with_suffix('.json').write_text(json.dumps(result, indent=2)+'\n')
    print(json.dumps({k:v for k,v in result.items() if k != 'records'}))


\end{lstlisting}

\subsection{Raw formal eta-product check}
Save as \path{check_raw_formal.py}.
\begin{lstlisting}
from math import comb
import numpy as np
import json
from pathlib import Path


def mul(a,b,N,mod):
 return np.convolve(a,b)[:N+1]%mod

def eta_product(N, exponents, mod):
 a=np.zeros(N+1,dtype=np.int64);a[0]=1
 for d,e in exponents.items():
  for n in range(1,N//d+1):
   step=d*n
   b=np.zeros(N+1,dtype=np.int64)
   for k in range(N//step+1):
    if e>=0:
     if k>e: break
     coeff=(-1)**k*comb(e,k)
    else: coeff=comb(-e+k-1,k)
    b[k*step]=coeff%mod
   a=mul(a,b,N,mod)
 return a

def substitute_p(a,p,N):
 b=np.zeros(N+1,dtype=np.int64)
 b[:p*(N//p)+1:p]=a[:N//p+1]
 return b

def check(p,s,m,N,extra=1):
 mod=p**(3*s+extra)
 assert (N+1)*(mod-1)**2<2**63
 R=eta_product(N,{1:12,4:12,2:-24},mod)
 theta=eta_product(N,{2:20,1:-8,4:-8},mod)
 rr=mul(R,R,N,mod)
 u=np.zeros(N+1,dtype=np.int64);u[1:]=rr[:N]
 curr=np.zeros(N+1,dtype=np.int64);prev=curr.copy();power=curr.copy();power[0]=1
 lim=min((m*p**s-1)//2,N)
 lower=(m*p**(s-1)-1)//2
 for j in range(lim+1):
  aj=pow(comb(2*j,j),3,mod)
  curr=(curr+aj*power)%mod
  if j<=lower:prev=(prev+aj*power)%mod
  if j<lim:power=mul(power,u,N,mod)
 defect=(mul(theta,substitute_p(prev,p,N),N,mod)-mul(substitute_p(theta,p,N),curr,N,mod))%mod
 req=p**(3*s)
 bad=[(int(j),int(v)) for j,v in enumerate(defect) if v%req]
 nonzero=[int(j) for j,v in enumerate(defect) if v]
 return dict(p=p,s=s,m=m,q_even_degree=N,required_exponent=3*s,extra_precision=extra,failures=bad[:5],first_nonzero_at_extra_precision=nonzero[:8])

if __name__=='__main__':
 out=[]
 for p,s,m,N in [(3,1,1,80),(3,2,1,100),(3,3,1,150),(3,4,1,180),(5,1,1,100),(5,2,1,150),(5,3,1,180),(5,3,2,220),(5,3,3,260),(5,3,5,360),(7,2,1,180),(7,2,3,220),(11,2,1,200)]:
  r=check(p,s,m,N)
  out.append(r);print(json.dumps(r),flush=True)
 Path(__file__).with_suffix('.json').write_text(json.dumps(out,indent=2))

\end{lstlisting}

\subsection{Supersingular half-cutoff recurrence}
Save as \path{check_supersingular.py}.
\begin{lstlisting}
"""Exact finite checks of the half-cutoff supersingular recurrence.

This program does not establish an untested prime or level.
"""
from pathlib import Path
import json
from math import comb

def strip(n, p):
    v = 0
    while n % p == 0:
        n //= p
        v += 1
    return n, v

def val(n, p, e):
    n %= p**e
    return {'value': e, 'exact': False} if not n else {'value': strip(n,p)[1], 'exact': True}

def levels(p, max_s):
    e=3*max_s+5
    mod=p**e
    ends={(p**s-1)//2:s for s in range(max_s+1)}
    out={}
    unit=1;v=0;total=0
    for k in range((p**max_s+1)//2):
        term=0 if 3*v>=e else pow(unit,3,mod)*pow(p,3*v,mod)%mod
        total=(total+term)%mod
        if k in ends:
            s=ends[k];out[s]=total
        a,va=strip(2*k+1,p);b,vb=strip(2*k+2,p)
        unit=unit*a*pow(b,-1,mod)%mod;v+=va-vb
        assert v>=0
    direct=sum(pow(comb(2*k,k),3,mod)*pow(64,-k,mod) for k in range((p+1)//2))%mod
    assert out[1]==direct
    result=[]
    for s in range(1,max_s+1):
        row={'p':p,'s':s,'precision':e,'T':out[s],'valuation_T':val(out[s],p,e)}
        if s>=2:
            defect=(out[s]-p*p*out[s-2])%mod
            row.update(defect=defect,valuation_defect=val(defect,p,e),target=3*s-1)
            row['proved_bound']=3*s-1-int(p==3)
            assert row['valuation_defect']['value']>=row['proved_bound']
        result.append(row)
    return result

if __name__=='__main__':
    rows=[]
    for p,s in [(3,10),(7,6),(11,5),(19,4),(23,4),(31,3),(43,3),(47,3)]:
        chunk=levels(p,s);rows.extend(chunk)
        print(json.dumps({'p':p,'T_values':[r['valuation_T']['value'] for r in chunk],
                          'defect_values':[r['valuation_defect']['value'] for r in chunk if r['s']>=2]}),flush=True)
    Path(__file__).with_suffix('.json').write_text(json.dumps({
        'scope':'Finite checks, not a proof of untested cases.',
        'recurrence_pairs':sum(r['s']>=2 for r in rows),
        'all_passed':True,'records':rows},indent=2)+'\n')

\end{lstlisting}

\begin{thebibliography}{99}
\bibitem{Swisher}
H. Swisher, \emph{On the supercongruence conjectures of van Hamme},
Research in the Mathematical Sciences \textbf{2} (2015), 18;
\href{https://doi.org/10.1186/s40687-015-0037-6}{doi:10.1186/s40687-015-0037-6};
\href{https://arxiv.org/abs/1504.01028}{arXiv:1504.01028}, Table~2, (H.3).

\bibitem{Guo} V.~J.~W.~Guo, \emph{Dwork-type $q$-congruences through
the $q$-Lucas theorem}, Revista de la Real Academia de Ciencias Exactas,
F\'isicas y Naturales. Serie A. Matem\'aticas \textbf{119} (2025), 97;
\href{https://doi.org/10.1007/s13398-025-01762-9}{doi:10.1007/s13398-025-01762-9};
\href{https://arxiv.org/abs/2310.15207}{arXiv:2310.15207}, equations (1.9)--(1.10).

\bibitem{Beukers}
F. Beukers, \emph{Supercongruences using modular forms},
\href{https://arxiv.org/abs/2403.03301v3}{arXiv:2403.03301v3},
especially the modular trace argument and Appendix~C, Section~9.1.

\bibitem{BV3} F.~Beukers and M.~Vlasenko, \emph{Dwork crystals III:
from excellent Frobenius lifts towards supercongruences},
\href{https://arxiv.org/abs/2105.14841}{arXiv:2105.14841}.

\bibitem{Katz}
N. M. Katz, \emph{$p$-adic properties of modular schemes and modular forms},
in Modular Functions of One Variable III, Lecture Notes in Mathematics
\textbf{350} (1973), 69--190;
\href{https://web.math.princeton.edu/~nmk/old/padicpropMFMS.pdf}{author's scan}.
See Proposition~2.7.1, Sections~3.1--3.4, and Appendix~A2, in particular
equations (A2.2.1) and (A2.2.7), printed pp.~175--177.

\bibitem{BM}
P. Berthelot and W. Messing, \emph{Th\'eorie de Dieudonn\'e cristalline I},
Ast\'erisque \textbf{63} (1979), 17--37;
\href{https://numdam.org/item/AST_1979__63__17_0.pdf}{original text}.
See Section~2.1 for connections and lift comparison, Theorem~2 and the
Frobenius--Verschiebung constructions in Section~2, and Section~3,
equation~(3.1), for the abelian-scheme cohomology and Hodge filtration.

\bibitem{Shirai}
R. Shirai, \emph{$q$-deformation with $(\varphi,\Gamma)$ structure of the
de Rham cohomology of the Legendre family of elliptic curves},
\href{https://arxiv.org/abs/2006.12310}{arXiv:2006.12310}.
The introduction records the Legendre basis, crystalline comparison,
cup product, and the sign $(-1)^{(p-1)/2}$ in the unit-root normalization;
Definition~3.1 records Frobenius--Hodge divisibility.

\bibitem{Stacks}
The Stacks Project Authors, \emph{The Stacks Project},
\href{https://stacks.math.columbia.edu/tag/07JH}{Proposition~60.17.4 (07JH)},
crystalline Taylor comparison;
\href{https://stacks.math.columbia.edu/tag/07MZ}{Remark~60.24.14 (07MZ)},
integral completed crystalline--de Rham comparison; and
\href{https://stacks.math.columbia.edu/tag/07N5}{Theorem~60.26.4 (07N5)},
the canonical integral Frobenius map. Accessed 10 November 2026.

\bibitem{Ligozat}
G. Ligozat, \emph{Courbes modulaires de genre 1},
M\'emoires de la Soci\'et\'e Math\'ematique de France \textbf{43} (1975),
5--80, Proposition~3.2.8;
\href{https://www.numdam.org/article/MSMF_1975__43__5_0.pdf}{original text}.

\bibitem{DLMF}
NIST Digital Library of Mathematical Functions,
\href{https://dlmf.nist.gov/20.7}{Section~20.7}, theta identities and
transformations. Accessed 10 November 2026.
\end{thebibliography}
\end{document}
